% {{Document title}} — skeleton from the research template.
% Compile: tectonic main.tex   (TeX Live 2023+ also works)
% Every number in this file comes from a script in scripts/, with its output in data/. Nothing is typed
% by hand. If you cannot name the script next to a number, the number does not go in.
\documentclass[11pt,a4paper]{article}
\usepackage[utf8]{inputenc}\usepackage[T1]{fontenc}
\usepackage{amsmath,amssymb,amsthm,microtype,xcolor,booktabs}
\usepackage[margin=2.5cm]{geometry}\linespread{1.06}\setlength{\emergencystretch}{3em}\hfuzz=2.5pt
\usepackage[colorlinks=true,linkcolor=blue!60!black,citecolor=blue!60!black]{hyperref}

\newtheorem{theorem}{Theorem}\newtheorem{proposition}[theorem]{Proposition}
\newtheorem{lemma}[theorem]{Lemma}\newtheorem{corollary}[theorem]{Corollary}
\theoremstyle{definition}\newtheorem{conjecture}[theorem]{Conjecture}
\newtheorem{hypothesis}[theorem]{Hypothesis}
\theoremstyle{remark}\newtheorem{remark}[theorem]{Remark}

% Markers. \todo is visible in the PDF on purpose: a draft should look like a draft.
\newcommand{\todo}[1]{\textcolor{red!70!black}{\small\sffamily[#1]}}
% Every number that comes from a script carries its provenance in the source.
\newcommand{\num}[2]{#1\footnote{Computed by \texttt{#2}; output in \texttt{data/}.}}

\title{{{Title}}}
\author{{{Author}}\thanks{{{Affiliation. Correspondence: ...}}}}
\date{{{date}} --- \todo{working draft, not for circulation until the readings of \S\ref{sec:status} are recorded}}

\begin{document}\maketitle

\begin{abstract}
\todo{The abstract is the most frequent site of drift between what is claimed and what is proved. After
every rewrite of a theorem, re-read this paragraph against the statement.}
{{What object, what is proved about it, under what hypothesis, and what is left open --- in that order.
The open part belongs in the abstract, not only in the last section.}}
\end{abstract}

\section{Introduction}\label{sec:intro}

\paragraph{The problem.} {{The object, defined, and why the question is the question. A reader from a
neighbouring field should get this far.}}

\paragraph{Results.} {{Each result in one sentence, with its label: proved; proved under a stated
hypothesis; conjectured with numerical support. The list must agree with \texttt{../KNOWLEDGE.md} \S2 and
with the theorem statements below --- check all three against each other, every revision.}}

\paragraph{What is not claimed.} {{The honest framing, stated here and not buried. What this does not
show, which standard question it does not touch, and what would be needed for the claim a reader might
otherwise infer. This paragraph is part of the result.}}

\paragraph{Related work.} {{What is known, what is ours, and what was checked against which source}}
(\texttt{../LITERATURE.md}), e.g.~\cite{Example2026}. {{A novelty claim without a named check does not
appear in this paper.}}

\section{Notation and set-up}\label{sec:setup}
{{Every object used later, defined once, in the normalisation the field uses. Where our normalisation
differs from a cited source's, say so here rather than in a footnote at the point of use.}}

\section{{{The main result}}}\label{sec:main}

\begin{theorem}\label{thm:main}
{{Self-contained: a reader who opens the paper at this page can tell what the object is and what the
hypotheses are. No forward references in a statement.}}
\end{theorem}

\begin{proof}
{{Structured in steps, each doing one thing. Every borrowed result is cited with its hypotheses checked
against the object here --- explicitly, in the text, not in the author's head. Every parameter introduced
is accompanied by its constraints, and the admissible set is shown non-empty.}}
\end{proof}

\section{Numerical verification}\label{sec:numerics}
{{What was computed, at what scale, with which script, and what the agreement is --- in units of an error
estimated two ways. A table of two independent routes for every constant. Say what the computation does
not and cannot distinguish: no finite computation separates a bounded quantity from one growing like
$\log\log$.}}

\begin{table}[htbp]\centering\small
\caption{{{What is compared.}} Regenerate with \texttt{python3 scripts/{{script}}.py}; output in
\texttt{data/{{file}}}.}
\begin{tabular}{@{}lrrr@{}}\toprule
{{quantity}} & route 1 & route 2 & digits agreeing\\\midrule
{{...}} & {{...}} & {{...}} & {{...}}\\
\bottomrule\end{tabular}\end{table}

\section{What remains}\label{sec:open}
{{The open points, stated exactly: what is missing, what strength of input would suffice, what the known
obstruction is, and what has been tried. ``We expect'' is not a statement; ``an estimate of the form X
with exponent below Y would give Z, and the best available is W'' is. Do not describe an open problem as
a lack of rigour, and do not call a technicality anything that has not been carried out.}}

\section{Status of this document}\label{sec:status}
\todo{Delete this section only when the document leaves the repository. Until then it is the honest
summary: which statements have been read by whom, and which have been read by nobody.}
{{Readings: who read what, at which revision. Which statements have had an adversarial reading by a reader
with no session context, and which have not. See \texttt{STATUS.md} for the dated record and
\texttt{../ERRATA.md} for every correction.}}

\section*{Disclosure}
\todo{Keep this section. Adapt it to what actually happened; do not soften it.}
{{How this work was produced: which tools, used how, by whom, with what checked by a human and what not.
State the class of error the process is known not to catch (see \texttt{../../METHOD.md}). If an AI system
was used, say so plainly, here, and not only in an acknowledgement.}}

\section*{Reproducibility}
Every number in this paper is produced by a script in \texttt{scripts/}, named at the point of use, with
its output in \texttt{data/}. {{The commands, in full, with runtimes.}} The state of the programme is
\texttt{../KNOWLEDGE.md} \S0b; every correction ever applied is \texttt{../ERRATA.md}.

\bibliographystyle{plain}
\bibliography{refs}
\end{document}
